\documentclass[12pt]{article}
\usepackage{amsmath,amssymb,amsfonts}
\usepackage[margin=1in]{geometry}
\begin{document}
\begin{center}
{\large\bfseries 8th Irish Mathematical Olympiad\\6 May 1995}
\end{center}
\bigskip\noindent\textbf{Paper 1}\bigskip
\begin{enumerate}
\item There are $n^2$ students in a class. Each week all the students participate in a table quiz. Their teacher arranges them into $n$ teams of $n$ players each. For as many weeks as possible, this arrangement is done in such a way that any pair of students who were members of the same team one week are not on the same team in subsequent weeks. Prove that after at most $n + 2$ weeks, it is necessary for some pair of students to have been members of the same team on at least two different weeks.

\item Determine, with proof, all those \textbf{integers} $a$ for which the equation
\[x^2 + axy + y^2 = 1\]
has infinitely many distinct \textbf{integer} solutions $x$, $y$.

\item Let $A$, $X$, $D$ be points on a line, with $X$ between $A$ and $D$. Let $B$ be a point in the plane such that $\angle ABX$ is greater than $120^\circ$, and let $C$ be a point on the line between $B$ and $X$. Prove the inequality
\[2|AD| \ge \sqrt{3}(|AB| + |BC| + |CD|).\]

\item Consider the following one-person game played on the $x$-axis. For each integer $k$, let $X_k$ be the point with coordinates $(k, 0)$. During the game discs are piled at some of the points $X_k$. To perform a move in the game, the player chooses a point $X_j$ at which at least two discs are piled and then takes two discs from the pile at $X_j$ and places one of them at $X_{j-1}$ and one at $X_{j+1}$.

To begin the game, $2n + 1$ discs are placed at $X_0$. The player then proceeds to perform moves in the game for as long as possible. Prove that after $n(n+1)(2n+1)/6$ moves no further moves are possible, and that, at this stage, one disc remains at each of the positions
\[X_{-n}, X_{-n+1}, \ldots, X_{-1}, X_0, X_1, \ldots, X_{n-1}, X_n.\]

\item Determine, with proof, all real-valued functions $f$ satisfying the equation
\[xf(x) - yf(y) = (x-y)f(x+y),\]
for all real numbers $x$, $y$.
\end{enumerate}

\newpage
\begin{center}
{\large\bfseries 8th Irish Mathematical Olympiad\\6 May 1995}
\end{center}
\bigskip\noindent\textbf{Paper 2}\bigskip
\begin{enumerate}
\setcounter{enumi}{5}
\item Prove the inequalities
\[n^n \le (n!)^2 \le [(n+1)(n+2)/6]^n,\]
for every positive integer $n$.

\item Suppose that $a$, $b$ and $c$ are complex numbers, and that all three roots $z$ of the equation
\[x^3 + ax^2 + bx + c = 0\]
satisfy $|z| = 1$ (where $|\;|$ denotes absolute value). Prove that all three roots $w$ of the equation
\[x^3 + |a|x^2 + |b|x + |c| = 0\]
also satisfy $|w| = 1$.

\item Let $S$ be the square consisting of all points $(x,y)$ in the plane with $0 \le x, y \le 1$. For each real number $t$ with $0 < t < 1$, let $C_t$ denote the set of all points $(x,y) \in S$ such that $(x,y)$ is on or above the line joining $(t,0)$ to $(0, 1-t)$.
Prove that the points common to all $C_t$ are those points in $S$ that are on or above the curve $\sqrt{x} + \sqrt{y} = 1$.

\item We are given three points $P, Q, R$ in the plane. It is known that there is a triangle $ABC$ such that $P$ is the mid-point of the side $BC$, $Q$ is the point on the side $CA$ with $|CQ|/|QA| = 2$, and $R$ is the point on the side $AB$ with $|AR|/|RB| = 2$. Determine, with proof, how the triangle $ABC$ may be constructed from $P$, $Q$, $R$.

\item For each integer $n$ such that $n = p_1p_2p_3p_4$, where $p_1, p_2, p_3, p_4$ are distinct primes, let
\[d_1 = 1 < d_2 < d_3 < \cdots < d_{15} < d_{16} = n\]
be the sixteen positive integers that divide $n$.
Prove that if $n < 1995$, then $d_9 - d_8 \ne 22$.
\end{enumerate}
\end{document}
